\begingroup
\shorthandoff{"<>} % Previene errores con babel y caracteres especiales en \texttt
\setlength{\arrayrulewidth}{1pt}
\arrayrulecolor{vlepsbased}
\begin{longtable}{|p{15cm}|}
  
  % --- CABECERA INICIAL ---
  \hline
  \rowcolor{maincolor} 
  \multicolumn{1}{|l|}{\textcolor{white}{\textbf{Caso de Uso}}} \\
  \hline
  \endfirsthead
  
  % --- CABECERA EN PÁGINAS SIGUIENTES ---
  \multicolumn{1}{c}{{\tablename\ \thetable{} -- Continuación}} \\
  \hline
  \rowcolor{maincolor} 
  \multicolumn{1}{|l|}{\textcolor{white}{\textbf{Caso de Uso (Continuación)}}} \\
  \hline
  \endhead
  
  % --- PIE DE PÁGINA (mientras no sea la última) ---
  \hline \multicolumn{1}{|r|}{\textit{Continúa en la siguiente página...}} \\ \hline
  \endfoot
  
  % --- PIE FINAL (Donde aparece el nombre ahora) ---
  \hline
  \caption{Caso de Uso CU1 — Interacción con el Tutor Socrático.}
  \label{TB:CU1_TUTOR_SOCRATICO_V3}
  \endlastfoot

  % Fila 1: ID y Actor
  \rowcolor{ulepsbased}
  \begin{tabularx}{\linewidth}{@{}p{7.5cm}p{7.5cm}@{}}
    \textbf{Identificador:} CU1 & \textbf{Actor principal:} Alumno 
  \end{tabularx} \\ \hline

  % Fila 2: Nombre
  \textbf{Nombre:} Interacción con el tutor socrático \\ \hline

  % Fila 3: Descripción
  \rowcolor{ulepsbased}
  \textbf{Descripción:} \newline 
  El alumno envía una pregunta al tutor socrático, opcionalmente acompañada del código de su editor y la última salida de ejecución. El sistema modera la entrada, genera una respuesta pedagógica en streaming (SSE) y opcionalmente modera la salida en paralelo. \\ \hline

  % Fila 4: Precondiciones
  \textbf{Precondiciones:} \newline
  -- El alumno ha iniciado sesión en la plataforma. \newline
  -- El servicio Ollama está activo y los modelos (principal y moderador) configurados están disponibles. \newline
  -- Existe una sesión de chat activa o se creará una nueva automáticamente. \\ \hline

  % Fila 5: Flujo principal
  \rowcolor{ulepsbased}
  \textbf{Flujo principal:} \newline
  1. El alumno escribe un mensaje en el chat y pulsa enviar.\newline
  2. El sistema valida que el mensaje no exceda la longitud permitida (2000 caracteres).\newline
  3. El sistema obtiene o crea la sesión de chat si esta no existe.\newline
  4. El sistema guarda el mensaje del alumno en la base de datos.\newline
  5. El sistema comprueba si es el primer mensaje de usuario y si es el caso asigna los primeros 50 caracteres como título. Actualiza la fecha de la última modificación.\newline
  6. El sistema lee la configuración global para determinar los modelos a utilizar, los modos de moderación y el umbral de palabras para la moderación de output.\newline
  7. Si la moderación de input está activada envía el texto del alumno al LLM de moderación, si el veredicto es OK, continúa.\newline
  8. El sistema construye el array de mensajes para Ollama: \texttt{SYSTEM\_PROMPT} + contexto de código (si existe) + últimos 10 mensajes del historial.\newline
  9. El sistema inicia el streaming asíncrono.\newline
  10. Por cada token recibido, el sistema lo transmite en tiempo real hacia la interfaz del usuario vía SSE.\newline
  11. Si la moderación de salida está activada, el sistema evalúa periódicamente el texto acumulado de forma asíncrona.\newline
  12. La interfaz de usuario recibe los fragmentos de manera progresiva y los visualiza de forma continua en el chat.\newline
  13. Al finalizar la generación de la respuesta, el sistema notifica el cierre de la transmisión y guarda la respuesta completa del asistente en la base de datos.\newline
  14. La interfaz de usuario actualiza el estado de la sesión activa y refresca el historial de conversaciones en la barra lateral. \\ \hline

  % Fila 6: Flujos alternativos
  \textbf{Flujos alternativos:} \newline
  \textbf{FA-01:} Moderación de input rechaza el mensaje (Paso 7):  \newline
  1. El sistema marca el mensaje como moderado en la base de datos. \newline
  2. Se envía respuesta predeterminada y no se invoca al LLM tutor. \newline
  \textbf{FA-02:} Moderación de output rechaza un fragmento (Paso 11):  \newline
  1. El sistema marca la respuesta como moderada en la base de datos. \newline
  2. Se interrumpe la generación y se envía respuesta predeterminada. \\ \hline

  % Fila 7: Postcondiciones
  \rowcolor{ulepsbased}
  \textbf{Postcondiciones:} \newline
  1. El mensaje del alumno se almacena en la base de datos. \newline
  2. La respuesta del tutor se almacena en la base de datos. \newline
  3. La sesión se actualiza o renombra si es necesario. \\ \hline

\end{longtable}
\endgroup